\documentclass[11pt,a4paper]{article}

\usepackage[margin=1.15in]{geometry}
\usepackage{amsmath,amssymb,amsthm,mathtools}
\usepackage{microtype}
\usepackage{enumitem}
\usepackage{hyperref}
\usepackage[T1]{fontenc}
\usepackage{lmodern}
\usepackage{setspace}
\setstretch{1.12}

\hypersetup{colorlinks=true,linkcolor=black,urlcolor=blue}

\newcommand{\Q}{\mathbb{Q}}
\newcommand{\Z}{\mathbb{Z}}
\newcommand{\C}{\mathbb{C}}
\newcommand{\Pspace}{\mathbb{P}}

\newtheorem{theorem}{Theorem}
\newtheorem{lemma}[theorem]{Lemma}
\theoremstyle{definition}
\newtheorem{definition}[theorem]{Definition}
\theoremstyle{remark}
\newtheorem{remark}[theorem]{Remark}

\title{\textbf{The Cycle Class Map, the Reverse Arrow,\\[0.25em]
and a Classical Constructor on Three Fourfolds}}
\author{Benjamin Stanley Frohman}
\date{20 September 2026}

\begin{document}
\maketitle

\begin{abstract}
The cycle class map
\(\mathrm{cl}:Z^k(X)_{\Q}\to H^{2k}(X,\Q)\)
is constructed from the fundamental class and Poincar\'e duality. Its
image consists of Hodge classes. A section of \(\mathrm{cl}\) over
\(\mathrm{Hdg}^k(X)\) is the Hodge conjecture. That section is written
explicitly on the Klein quadric \(Q^4\subset\Pspace^5\), on
\(\Pspace^4\), and on \(\Pspace^2\times\Pspace^2\), with equations,
Schubert identities, and intersection formulae. The same rule does not
extend to a general fourfold. The remaining emptiness is the Hodge
conjecture.
\end{abstract}

\section{The easy arrow}

The last geometric object on screen is not a mystery class. It is the
cycle class of a subvariety: the one arrow that is already a theorem.
The line ``Decoding still required'' refers to the reverse arrow, which
remains open on a general variety.

In the film a closed algebraic subset \(Z\subset X\) of complex
codimension \(k\) is named, and then its class
\[
[Z]\in H^{2k}(X,\Q)\cap H^{k,k}(X).
\]
That symbol \([Z]\) is the cycle class of \(Z\). Every later frame asks
whether an arbitrary Hodge class is a rational combination of objects
of this same kind.

The yellow loop drawn on the wireframe is a picture, not a construction.
A genuine algebraic subvariety of complex codimension \(k\) is cut out
by polynomial equations and has complex dimension \(\dim X-k\). On an
elliptic curve the algebraic classes in positive codimension are points,
not real circles. The animation is showing the role of \(Z\), not a
particular equation.

\section{How the class is produced}

Let \(X\) be a smooth complex projective variety of dimension \(n\), and
let \(Z\subset X\) be a closed irreducible subvariety of codimension
\(k\). The complex structure orients the regular locus of \(Z\) as a
real \(2(n-k)\)-dimensional manifold. That orientation supplies a
fundamental class in homology,
\[
[Z]_{\mathrm{hom}}\in H_{2n-2k}(X,\Z).
\]
Poincar\'e duality on the compact oriented manifold \(X\) converts it
into a cohomology class
\[
[Z]\;:=\;\mathrm{PD}\bigl([Z]_{\mathrm{hom}}\bigr)
\in H^{2k}(X,\Z)\subset H^{2k}(X,\Q).
\]
Equivalently, \(Z\) defines a closed current of integration: for a
smooth form \(\varphi\) of degree \(2n-2k\),
\[
T_Z(\varphi)\;:=\;\int_{Z^{\mathrm{reg}}}\varphi.
\]
The current \(T_Z\) is closed, of type \((k,k)\), and represents the
same class. Extending by linearity from irreducible subvarieties to
rational algebraic cycles gives the cycle class map
\[
\mathrm{cl}\colon\quad Z^k(X)_{\Q}\;\longrightarrow\; H^{2k}(X,\Q).
\]
The image of \(\mathrm{cl}\) is the space of algebraic classes.

\section{Why the class is automatically a Hodge class}

A current of type \((k,k)\) pairs to zero against every form of type
\((p,q)\) with \((p,q)\neq(n-k,n-k)\). Passing to cohomology, the class
therefore has vanishing off-diagonal Hodge components:
\[
[Z]\in H^{k,k}(X).
\]
Combined with rationality,
\[
[Z]\in H^{2k}(X,\Q)\cap H^{k,k}(X).
\]
That intersection is the definition of a Hodge class. This is the easy
arrow:
\[
\text{geometry }Z\;\longrightarrow\;\text{code }[Z].
\]
It holds for every algebraic cycle, with no restriction on \(k\). The
line ``Geometry \(\to\) code. Always. No exceptions.'' is this theorem.

In codimension one the same class has a second, more explicit name. If
\(Z\) is a divisor, then
\[
[Z]\;=\;c_1\bigl(\mathcal{O}_X(Z)\bigr),
\]
the first Chern class of the associated line bundle. The exponential
sequence produces every integral \((1,1)\) class in this way, and GAGA
promotes the holomorphic bundle to an algebraic divisor. That is why
\(k=1\) is settled. There is no analogous exponential sequence for
\(k\ge 2\).

\section{What decode would mean}

The cycle class map always runs in one direction. Decoding a Hodge class
would mean inverting \(\mathrm{cl}\) on the Hodge subspace: given
\[
\gamma\in H^{2k}(X,\Q)\cap H^{k,k}(X),
\]
produce finitely many subvarieties \(Z_i\subset X\) of codimension \(k\)
and rationals \(a_i\in\Q\) such that
\[
\gamma\;=\;\sum_i a_i[Z_i].
\]
That is the Hodge conjecture. The map \(\mathrm{cl}\) is known. A
section of \(\mathrm{cl}\) over the Hodge classes is not known, except
in the cases
\begin{itemize}[leftmargin=1.6em]
\item \(k=1\), by Lefschetz \((1,1)\);
\item \(k=\dim X-1\), by hard Lefschetz applied to the previous case;
\item \(\dim X\le 3\), because those two degrees exhaust the interesting
even cohomology.
\end{itemize}
The first open geometric case is a codimension-\(2\) class on a
fourfold. The class \([Z]\) can be written down whenever \(Z\) is given,
and cannot in general be reconstructed from \(\gamma\) alone.

\section{The first open case}

Let \(X\) be a general fourfold and let \(\gamma\in\mathrm{Hdg}^2(X)\)
be arbitrary. The missing object is a finite collection of surfaces
\(Z_i\subset X\) and rationals \(a_i\in\Q\) such that
\[
\gamma\;=\;\sum_i a_i[Z_i].
\]
That equality is the Hodge conjecture in its first open case. Writing
the \(Z_i\) down from \(\gamma\) would be a section of \(\mathrm{cl}\).
The cycle class map itself is constructed as in the preceding sections.

The two slots are therefore these. A rule \(\gamma\mapsto(Z_i,a_i)\) is
a section of
\[
\mathrm{cl}\colon\mathrm{CH}^2(X)_{\Q}\longrightarrow\mathrm{Hdg}^2(X)
\]
on a fourfold. A fourfold \(X\) together with a class
\(\gamma\in\mathrm{Hdg}^2(X)\setminus\mathrm{im}(\mathrm{cl})\) is a
counterexample. The Hodge conjecture is the claim that the first object
exists for every smooth complex projective \(X\), and that the second
object therefore does not.

The blanks below are filled only on fourfolds where the surfaces are
classical and the identities can be checked. They cannot be filled for
a general fourfold. What follows is not a solution of Clay's problem.

\section{Input: the Klein quadric}

Let \(X=Q\) be the Klein quadric in \(\Pspace^5\), with Pl\"ucker
coordinates \((p_{01}:p_{02}:p_{03}:p_{12}:p_{13}:p_{23})\) and equation
\begin{equation}
\label{eq:plucker}
p_{01}\,p_{23}-p_{02}\,p_{13}+p_{03}\,p_{12}\;=\;0.
\end{equation}
This hypersurface is smooth of dimension \(4\). It is
\(\mathrm{Gr}(2,4)\), the Grassmannian of lines in \(\Pspace^3\).
Lefschetz hyperplane together with the topology of quadrics gives
\[
\dim_{\Q}H^4(Q,\Q)\;=\;2,\qquad H^{3,1}(Q)=H^{1,3}(Q)=0,
\]
so every class in \(H^4(Q,\Q)\) is of type \((2,2)\):
\[
\mathrm{Hdg}^2(Q)\;=\;H^4(Q,\Q).
\]
The Hodge conjecture is a theorem for quadrics. Equivalently,
\(\mathrm{Gr}(2,4)\) is cellular, its cohomology is spanned by Schubert
cycles, and every class is algebraic.

\section{Output: two explicit planes}

Write \(\Pspace^3\) with coordinates \([z_0:z_1:z_2:z_3]\).

\paragraph{\(\sigma\)-plane \(\Pi\).}
Lines through the point \([1:0:0:0]\):
\begin{equation}
\label{eq:sigma}
\Pi\;=\;\bigl\{\,p_{12}=p_{13}=p_{23}=0\,\bigr\}\;\subset\;Q.
\end{equation}
This is a linearly embedded \(\Pspace^2\subset Q\). In Schubert notation
it is the cycle of class \(\sigma_2\).

\paragraph{\(\rho\)-plane \(\Pi'\).}
Lines contained in the plane \(z_0=0\):
\begin{equation}
\label{eq:rho}
\Pi'\;=\;\bigl\{\,p_{01}=p_{02}=p_{03}=0\,\bigr\}\;\subset\;Q.
\end{equation}
This is a second linearly embedded \(\Pspace^2\subset Q\), of the opposite
family. In Schubert notation it is the cycle of class \(\sigma_{1,1}\).

Let \(h=c_1(\mathcal{O}_Q(1))\). Schubert calculus on \(\mathrm{Gr}(2,4)\)
gives the identity of classes
\begin{equation}
\label{eq:h2}
h^2\;=\;[\Pi]+[\Pi']
\end{equation}
in \(H^4(Q,\Z)\), which is the relation \(\sigma_1^2=\sigma_2+\sigma_{1,1}\).
The two classes \([\Pi]\) and \([\Pi']\) are linearly independent over
\(\Q\) and therefore a basis of \(\mathrm{Hdg}^2(Q)\). Their intersection
numbers are
\begin{equation}
\label{eq:int}
[\Pi]\cdot[\Pi]\;=\;1,\qquad
[\Pi']\cdot[\Pi']\;=\;1,\qquad
[\Pi]\cdot[\Pi']\;=\;0.
\end{equation}

\section{The rule, on this fourfold}

\textbf{Input.} A class \(\gamma\in\mathrm{Hdg}^2(Q)\).

\textbf{Expansion.} On the basis \eqref{eq:h2}--\eqref{eq:int},
\begin{equation}
\label{eq:expand}
\gamma\;=\;a[\Pi]+b[\Pi'],\qquad a,b\in\Q,
\end{equation}
where the coefficients are the intersection numbers
\[
a\;=\;\gamma\cdot[\Pi],\qquad b\;=\;\gamma\cdot[\Pi'].
\]

\textbf{Output.} The surfaces \(\Pi,\Pi'\) written in
\eqref{eq:sigma}--\eqref{eq:rho} and the coefficients \(a,b\).

\begin{lemma}
\(\gamma=a[\Pi]+b[\Pi']\).
\end{lemma}

This is the expansion of \(\gamma\) in a basis of algebraic classes.

\begin{lemma}
The same procedure works for every \(\gamma\in\mathrm{Hdg}^2(Q)\).
\end{lemma}

Because \(\mathrm{Hdg}^2(Q)=\Q[\Pi]\oplus\Q[\Pi']\).

\begin{theorem}
\label{thm:quadric}
Let \(Q\subset\Pspace^5\) be the Klein quadric \eqref{eq:plucker}. The
assignment
\[
\gamma\;\longmapsto\;\bigl(\Pi,\,\Pi',\,\gamma\cdot[\Pi],\,\gamma\cdot[\Pi']\bigr)
\]
is a section of
\(\mathrm{cl}:Z^2(Q)_{\Q}\to\mathrm{Hdg}^2(Q)\).
\end{theorem}

\section{Two further fourfolds}

If one wants a single surface rather than two families: take
\(X=\Pspace^4\). Then \(\mathrm{Hdg}^2(\Pspace^4)=\Q\,h^2\), and \(h^2\)
is the class of any linearly embedded plane
\begin{equation}
\label{eq:p4}
Z\;=\;\bigl\{x_3=x_4=0\bigr\}\;\subset\;\Pspace^4.
\end{equation}
The rule is \(\gamma=a[Z]\) with \(a\in\Q\). Every Hodge class on this
\(X\) is hit.

On \(\Pspace^2\times\Pspace^2\) the same pattern gives three surfaces:
\(\{\mathrm{pt}\}\times\Pspace^2\), \(\Pspace^2\times\{\mathrm{pt}\}\),
and a line-times-line \(\Pspace^1\times\Pspace^1\), with classes
\(h_1^2\), \(h_2^2\), and \(h_1 h_2\). These three classes are a basis of
\[
H^4(\Pspace^2\times\Pspace^2,\Q)=\mathrm{Hdg}^2(\Pspace^2\times\Pspace^2).
\]
The rule is the expansion of \(\gamma\) on that basis.

\section{What this does not fill}

The remaining demand was that the same procedure work for every
\(\gamma\) on a general fourfold. The rule above uses that
\(H^4(Q,\Q)\) is spanned by planes. That fails as soon as the fourfold
has Hodge classes that are not in the span of those planes.

Two facts must be kept separate.

\begin{enumerate}[leftmargin=1.6em]
\item A general cubic fourfold \(Y\subset\Pspace^5\) has
\(h^{3,1}(Y)=1\) and \(h^{2,2}(Y)=21\). The two planes of the Klein
quadric are not present, and \(H^4(Y,\Q)\) is not spanned by planes.
The Hodge conjecture in codimension two is nevertheless a theorem for
cubic fourfolds: Voisin, using the Fano variety of lines and the
incidence correspondence. For a \emph{very general} cubic one has
moreover \(\mathrm{Hdg}^2(Y)=\Q\,h^2\) by the Noether--Lefschetz
theorem in this degree, so the only Hodge class is the hyperplane
section of a cubic surface.
\item A very general hypersurface of degree \(d\ge 6\) in \(\Pspace^5\)
again has \(\mathrm{Hdg}^2=\Q\,h^2\) by Noether--Lefschetz, so the
Hodge conjecture in codimension two holds for that trivial reason.
The open content of the conjecture on fourfolds is the existence of a
constructor for those fourfolds that do carry extra rational Hodge
classes --- special hypersurfaces, Calabi--Yau fourfolds with
\(h^{3,1}>0\) and \(\mathrm{rk}\,\mathrm{Hdg}^2>1\), abelian fourfolds
beyond the known cases, and products not built from curves and
surfaces.
\end{enumerate}

On those varieties the constructor is a different geometric object
(Fano lines for cubics; nothing uniform beyond that). The blanks are
filled for \(Q^4\), for \(\Pspace^4\), and for
\(\Pspace^2\times\Pspace^2\). They remain empty for a general fourfold
in the sense of an arbitrary smooth complex projective fourfold. That
remaining emptiness is the Hodge conjecture.

\end{document}
